% -----------------------------------------------------------------------------
%  Figure 3.4 — Architecture de l'infrastructure d'hébergement en nuage
%  effectivement déployée. Chemin principal en haut, dérivations en dessous.
%  Les précisions de chaque lien sont portées par la note, et non par des
%  étiquettes glissées entre les blocs, où elles ne tiendraient pas.
% -----------------------------------------------------------------------------
\begin{tikzpicture}[x=1cm, y=1cm]

  % ===================== CHEMIN PRINCIPAL ====================================
  \node[blocfort, text width=2.3cm, anchor=north, minimum height=13mm] (vis)
       at (-6.2,0) {Visiteur\\ \textit{navigateur, téléphone}};

  \node[bloc, text width=2.9cm, anchor=north, minimum height=13mm] (r53)
       at (-2.7,0) {\textbf{Route 53}\\ zone hébergée\\
                    enregistrements A, AAAA et joker};

  \node[bloc, text width=3.3cm, anchor=north, minimum height=13mm] (cf)
       at (1.3,0) {\textbf{CloudFront}\\ réseau de diffusion\\
                   certificat \textbf{ACM}};

  \node[bloc, text width=2.9cm, anchor=north, minimum height=13mm] (s3)
       at (5.4,0) {\textbf{Amazon S3}\\ privé, chiffré, versionné};

  \draw[fleche] (vis.east) -- (r53.west);
  \draw[fleche] (r53.east) -- (cf.west);
  \draw[fleche] (cf.east) -- (s3.west);

  % ===================== RÉGLAGES ET DÉRIVATIONS =============================
  \node[bloc, text width=3.4cm, anchor=north, minimum height=13mm] (sup)
       at (-4.4,-2.5) {\textbf{Supabase}\\ base de données,
                       authentification, fonctions serveur};
  \draw[fleche] (vis.south) -- (vis.south |- sup.north)
        node[etiquette, pos=0.45, anchor=west, xshift=3pt]
        {interrogé directement par le navigateur};

  \node[blocgris, text width=3.3cm, anchor=north, minimum height=13mm] (regl)
       at (1.3,-2.5) {repli applicatif sur 403 et 404,\\
                      durées de cache différenciées,\\
                      en-têtes de sécurité};
  \draw[lien, dashed] (cf.south) -- (regl.north);

  \node[blocext, text width=3.0cm, anchor=north, minimum height=13mm] (gpu)
       at (5.4,-2.5) {\textbf{EC2 à processeur graphique}\\
                      photogrammétrie, 6 h, arrêtée hors usage};
  \draw[flechefine, draw=keyceorange] (s3.south) -- (gpu.north)
        node[etiquette, pos=0.5, anchor=west, xshift=3pt, text=keyceorange]
        {ponctuel};

  % ===================== ENCADRÉ TERRAFORM ===================================
  \begin{scope}[on background layer]
    \node[couche, fit=(r53)(cf)(s3)(regl)(gpu), draw=keycemid] (btf) {};
  \end{scope}
  \node[titrecouche, anchor=west] at ([xshift=6pt]btf.north west)
       {29 RESSOURCES DÉCRITES PAR TERRAFORM \textnormal{\itshape(19 créées à la mise en ligne)}};

  % ===================== NOTE ================================================
  \node[blocnote, anchor=north west, text width=14.2cm] at (-7.4,-4.9)
       {Le visiteur atteint le site en HTTPS ; Route 53 résout le nom, y compris
        le sous-domaine de chaque organisation grâce à l'enregistrement joker ;
        CloudFront sert le contenu et accède seul au bucket, par contrôle
        d'accès à l'origine. Aucun serveur ne fonctionne en permanence :
        l'application est un site statique. C'est ce choix d'architecture, et
        non une optimisation postérieure, qui ramène le coût récurrent
        d'hébergement à 730 FCFA par mois.};

\end{tikzpicture}
